%!TEX root = ../main.tex

\section{Discussion}\label{sec:discussion}

%=============================================================================
\subsection{Theory Assessment}
%=============================================================================

The endogenous premium model provides a coherent explanation for Strategy's
persistent NAV premium.  The core insight---that the premium is the present
value of an accumulation option---resolves the apparent anomaly: unlike a
closed-end fund discount driven by sentiment or agency costs, the premium
reflects genuine economic value created by the ability to issue equity above
NAV and purchase BTC accretively.

The model's predictions align well with calibrated data across six
dimensions (Section~\ref{sec:theory_validation}): positive premium from
accumulation capability, premium below steady-state implying positive carry,
financing channel selection matching the premium regime, dampened reflexivity
from premium compression, adequate survival margin, and procyclical leverage.

The fixed-point equation~\eqref{eq:pi_star_closed} provides a quantitative
tool for premium assessment.  At steady state, the fair premium is determined
by four fundamentals: BTC volatility $\sigma_S$ (positive vega), issuance
intensity $\alpha_N$ (moneyness), leverage $\ell$ (negative effect), and
discount rate $r$ (negative rho).  Changes in any of these fundamentals
shift the equilibrium, with the comparative statics providing directional
guidance.

%=============================================================================
\subsection{Where Is Strategy Currently?}
%=============================================================================

The indicator dashboard places Strategy in the ``MSTR dominates'' quadrant
of the regime map (Corollary~\ref{cor:regime_map}): IBGR is positive
(17.6\%) and PMRI is negative ($-0.71$).  The accumulation option is in the
money, the premium is underpriced relative to its long-run equilibrium,
and leverage is moderate ($\mathrm{ILE} = 1.49 < 2$).

However, several factors warrant caution:
\begin{itemize}
  \item The current premium ($\pi_0 = 0.242$) is below the ATM accretion
    threshold ($\pi^*_{\mathrm{ATM}} = 0.397$), meaning common equity
    issuance is currently mildly dilutive.  Strategy must rely on
    preferred stock and convertible debt to fund BTC purchases---channels
    with effective costs of 10--12.5\%.
  \item The preferred dividend burden of \$977M/year is a fixed obligation
    regardless of BTC price.  While the dividend coverage ratio is
    comfortable at 49$\times$, a sustained BTC decline to \$15,000--\$20,000
    would bring the leverage spiral condition
    (Corollary~\ref{cor:leverage_spiral}) into play.
  \item The 91.9\% three-year survival probability, while reassuring,
    implies an 8.1\% probability of NAV breach---non-trivial for a
    single-asset vehicle.
\end{itemize}

The critical near-term question is whether the premium recovers above the
ATM threshold, re-enabling the cheapest financing channel.  If BTC
appreciates significantly and the premium expands, the flywheel can
re-engage.  If the premium continues to compress, Strategy faces increasing
reliance on expensive preferred financing, degrading the capital structure
toward the tipping point of Theorem~\ref{thm:three_equilibria}.

%=============================================================================
\subsection{Implications for Investors}
%=============================================================================

\paragraph{MSTR vs.\ ETF.}
At the current calibration point, the indicators favour MSTR over a spot
BTC ETF for investors with sufficient risk tolerance:
\begin{itemize}
  \item Positive carry from $\mathrm{PMRI} < 0$ (expected premium
    expansion).
  \item Active accumulation option (IBGR $> 0$).
  \item Moderate leverage providing amplified upside.
\end{itemize}
However, the ETF is preferable for investors seeking pure BTC exposure
without the embedded risks of the capital structure, particularly given
the negative TEE (equity moving inversely to BTC at the margin).

\paragraph{Signal monitoring.}
The regime map provides actionable signals:
\begin{itemize}
  \item \textbf{PMRI crossing zero} signals the transition from
    positive to negative carry---the premium is now overpriced relative
    to fundamentals.
  \item \textbf{IBGR approaching zero} signals the flywheel is stalling
    ---accumulation is no longer accretive.
  \item \textbf{ILE exceeding 2} signals excessive leverage and
    elevated tail risk.
  \item \textbf{DCR falling below 10$\times$} signals preferred dividend
    sustainability concerns.
\end{itemize}

%=============================================================================
\subsection{Optimal Capital Structure Prescriptions}
%=============================================================================

The analysis yields concrete prescriptions for Strategy's financing:
\begin{enumerate}
  \item \textbf{Resume ATM when $\pi_t > 0.40$:} This is the accretion
    threshold at current leverage.  ATM issuance is the cheapest channel
    and should be used aggressively above this level.
  \item \textbf{Limit preferred issuance:} Each additional preferred
    series increases the dividend burden and the ATM accretion threshold
    ($\pi^*_{\mathrm{ATM}} = \log(\mathrm{ILE})$ rises with
    $P^{\mathrm{liq}}$).  The current \$977M annual dividend is already
    a material fixed cost.
  \item \textbf{Maintain cash buffer:} The \$2.25B cash reserve provides
    2.5 years of preferred dividend coverage in a zero-BTC-liquidation
    scenario.  This buffer is essential for surviving temporary drawdowns
    without forced sales.
  \item \textbf{Minimise WACBA:} Financing decisions should target the
    lowest WACBA given premium and leverage constraints, shifting
    channels as $\pi_t$ fluctuates.
\end{enumerate}

%=============================================================================
\subsection{Systemic Risk}
%=============================================================================

Strategy's success has catalysed corporate BTC adoption: over 70 public
companies now hold Bitcoin on balance sheet.  The market impact analysis
(Section~\ref{sec:market_impact_theory}) highlights two systemic
considerations:

\begin{enumerate}
  \item \textbf{Concentration risk:} Strategy holds 3.84\% of circulating
    BTC supply.  A forced liquidation at the full-stack distress threshold
    (\$24,200) would involve selling up to 761,068~BTC, roughly 30 days
    of average trading volume.  The price impact of such a sale, using
    Kyle lambda estimates, would be devastating.
  \item \textbf{Contagion:} If aggregate corporate holdings reach
    10--15\% of circulating supply, the contagion cascade mechanism
    (Theorem~\ref{thm:contagion}) becomes material.  A BTC decline that
    triggers one firm's distress threshold could propagate through
    correlated forced sales.
\end{enumerate}

The asymptotic stability result (Corollary~\ref{cor:asymptotic}) provides
a counterpoint: as BTC market depth grows ($2g_S + g_V > g_E$), the
reflexive feedback naturally weakens.  The question is whether market
maturation outpaces corporate adoption.

%=============================================================================
\subsection{Limitations and Future Work}
%=============================================================================

\paragraph{Limitations.}
\begin{enumerate}
  \item \emph{Physical measure:} We work under $\mathbb{P}$ throughout.
    A risk-neutral ($\mathbb{Q}$) formulation would enable derivative
    pricing and arbitrage-free valuation.
  \item \emph{Constant parameters:} The OU calibration assumes
    time-invariant $\kappa$, $\theta$, $\sigma_\pi$.  Regime-switching
    or time-varying parameters may better capture the structural
    break around the Q4~2024 premium spike.
  \item \emph{Single-firm focus:} Calibration is specific to Strategy.
    Generalisation to other BTC treasury vehicles requires
    firm-specific parameter estimation.
  \item \emph{Preferred stock modelling:} We treat preferred stock as a
    fixed obligation (liquidation value plus dividends).  The embedded
    call provisions, conversion optionality (STRK), and variable-rate
    features (STRC) introduce additional complexity not fully captured.
  \item \emph{Market impact estimation:} The Kyle lambda parameter
    $\eta$ is not directly estimated from our data.  We rely on
    literature estimates for illustrative analysis.
\end{enumerate}

\paragraph{Future work.}
Natural extensions include: (i)~risk-neutral pricing of the accumulation
option using derivatives data; (ii)~regime-switching models for $\kappa$
and $\theta$; (iii)~cross-sectional estimation across multiple BTC treasury
vehicles; (iv)~structural estimation of the feedback gain $G$ from
high-frequency issuance-price data; and (v)~welfare analysis of corporate
BTC adoption under the contagion cascade framework.
